\documentclass[11pt]{article}
\usepackage[a4paper,margin=27mm]{geometry}
\usepackage{amsmath,amssymb,amsthm}
\usepackage[T1]{fontenc}
\usepackage{lmodern}
\usepackage{microtype}
\usepackage[hidelinks]{hyperref}
\newtheorem{theorem}{Theorem}
\newtheorem{lemma}{Lemma}
\newtheorem{corollary}{Corollary}
\newcommand{\R}{\mathbb R}
\newcommand{\N}{\mathbb N}
\title{A graphon class without a finite-step representative\\
\large Counterexample to Conjecture 00000004408}
\author{}
\date{}

\begin{document}
\maketitle
\vspace{-1.8em}

\begin{abstract}
The symmetric graphon $W(x,y)=xy$ on the unit square has a nonatomic
distribution of function values. This distribution is unchanged by
measure-preserving relabellings and modification on null sets. Consequently,
no representative of its equivalence class is essentially countably valued,
and in particular no representative is a finite-step graphon. This refutes
the first universal assertion of Conjecture 00000004408. The accompanying
Lean 4 proof verifies the measure-theoretic obstruction, its invariance
under the generated equivalence relation, and the finite-step consequence.
\end{abstract}

\section{Statement and conventions}
The source defines $L^p$ reconstruction equivalence as equivalence of
graphons under measure-preserving transformations and zero-measure
modification. It then asserts that every class has a step representative
of minimal step count, that this count is a class invariant, and that for
purely jumping graphons it is the number of jumps.

The English and Chinese statements agree on these points. Neither states
a restriction to graphons that are already step functions. We use the
usual meaning of a finite-step graphon: there is a finite measurable
partition $[0,1]=A_1\sqcup\cdots\sqcup A_n$, up to null sets, and constants
$c_{ij}$ such that $U(x,y)=c_{ij}$ for almost every
$(x,y)\in A_i\times A_j$. Refining a partition is permitted. Our obstruction
also rules out a countable partition with one constant on each block
rectangle. An uncountable partition into singletons is not a finite- or
countable-step representation.

Write $\lambda$ for Lebesgue probability measure on $[0,1]$ and
$m=\lambda\otimes\lambda$. For a measurable, measure-preserving map
$\phi:[0,1]\to[0,1]$, let
\[
 U^\phi(x,y)=U(\phi(x),\phi(y)).
\]
We allow these relabellings without requiring invertibility, and take the
equivalence relation generated by such relabellings and equality $m$-almost
everywhere. Thus the argument covers, in particular, the convention using
only measure-preserving bijections. No limit of relabellings or unspecified
topological closure is used. The source's equivalence definition does not
depend on a numerical choice of $p$; approximation in an $L^p$ norm is a
different assertion.

\section{The counterexample}
Set
\[
 W:[0,1]^2\longrightarrow[0,1],\qquad W(x,y)=xy.
\]
This is a graphon: it is continuous, symmetric, and has values between
$0$ and $1$ on its stated domain.

\begin{lemma}\label{lem:countable}
For every countable set $S\subseteq\R$,
\[
 m\{(x,y): W(x,y)\in S\}=0.
\]
\end{lemma}
\begin{proof}
For each fixed $x\ne0$, the map $y\mapsto xy$ is injective, so the set
$\{y\in[0,1]:xy\in S\}$ is countable and has Lebesgue measure zero.
The only exceptional value $x=0$ is itself a null set. The inverse image
$W^{-1}(S)$ is measurable because $S$ is countable and $W$ is continuous.
Tonelli's theorem, applied to its indicator function, therefore gives
\[
 m(W^{-1}(S))
 =\int_{[0,1]}\lambda\{y\in[0,1]:xy\in S\}\,d\lambda(x)=0.
\]
The possible nonzero inner measure at $x=0$ has no effect on the integral.
\end{proof}

For any graphon $U$, define its value distribution by the pushforward
probability measure
\[
 \nu_U=U_*m,\qquad \nu_U(B)=m(U^{-1}(B))
 \quad\text{for every Borel set }B\subseteq\R.
\]

\begin{lemma}\label{lem:invariant}
The measure $\nu_U$ is invariant under the reconstruction equivalence
defined above.
\end{lemma}
\begin{proof}
If $U=V$ almost everywhere, their inverse images of every Borel set agree
up to a null set, hence $\nu_U=\nu_V$. If $V=U^\phi$, the map
$\phi\times\phi$ preserves the product measure $m$. Indeed, its pushforward
agrees with $m$ on measurable rectangles because
\[
 m\bigl((\phi\times\phi)^{-1}(A\times B)\bigr)
 =\lambda(\phi^{-1}(A))\lambda(\phi^{-1}(B))
 =\lambda(A)\lambda(B),
\]
and these rectangles determine the product probability measure. It follows
that
\[
 \nu_{U^\phi}=U_*\bigl((\phi\times\phi)_*m\bigr)=U_*m=\nu_U.
\]
Equality of measures is reflexive, symmetric, and transitive; invariance
therefore passes to the full equivalence relation generated by the two
operations, including operations traversed in reverse.
\end{proof}

\begin{theorem}\label{thm:main}
No graphon equivalent to $W(x,y)=xy$ is essentially countably valued.
In particular, the equivalence class of $W$ has no finite-step
representative.
\end{theorem}
\begin{proof}
Suppose $V$ is equivalent to $W$ and $V(x,y)\in S$ almost everywhere for a
countable set $S$. Then $\nu_V(S)=1$. Lemmas~\ref{lem:countable}
and~\ref{lem:invariant} give
\[
 1=\nu_V(S)=\nu_W(S)=0,
\]
a contradiction. A finite-step graphon with $n$ blocks is almost
everywhere valued in the finite set $\{c_{ij}:1\le i,j\le n\}$, so it is
excluded. The same reasoning applies to countably many blocks, since a
countable union of countable sets is countable.
\end{proof}

\begin{corollary}
The assertion that every reconstruction-equivalence class contains a
finite-step representative, hence one with a minimal finite step count,
is false. Therefore Conjecture 00000004408, as stated, is false.
\end{corollary}

This conclusion does not require deciding a convention for ``number of
jumps'' in two variables, or separately refuting an assertion limited to
graphons already known to have a step representation. It also does not
deny that general graphons can be approximated by step functions.

\section{Lean verification and the translation to graphons}
The project pins Lean \texttt{v4.19.0} and the following Mathlib commit:
\par\noindent
\texttt{c44e0c8ee63ca166450922a373c7409c5d26b00b}.
\par
The proof is in \texttt{lean/Main.lean}. The following mathematical steps
are checked by the Lean kernel:
\begin{enumerate}
\item \texttt{productKernel\_range}\par
      Together with the measurability and symmetry lemmas, this verifies
      the graphon conditions for $xy$.
\item \texttt{productKernel\_avoids\_countable}\par
      This proves Lemma~\ref{lem:countable}
      for every countable subset of $\R$, using product measure and
      nonatomic Lebesgue measure.
\item \texttt{equivalence\_preserves\_valueLaw}\par
      This proves Lemma~\ref{lem:invariant}, including its passage to
      \texttt{Relation.EqvGen}.
\item \texttt{no\_finite\_step\_in\_productKernel\_class}\par
      This proves the final nonexistence statement using the more general
      lemma excluding all essentially countably valued representatives.
\end{enumerate}

For the measure-space model, Lean uses $\R$ with the measure
$\lambda_0=\operatorname{volume}|_{[0,1]}$, then
$\lambda_0\otimes\lambda_0$. The file verifies that both are probability
measures. This is the unit-interval model with a null exterior, not a claim
that $xy$ has values in $[0,1]$ on all of $\R^2$. Kernels are real-valued
functions that are measurable almost everywhere for this product measure,
so null modifications are permitted.

The formal finite-step predicate allows an arbitrary map
$b:\R\to\operatorname{Fin}(n)$ and arbitrary weights $c_{ij}$, with
$U(x,y)=c_{b(x)b(y)}$ almost everywhere. Every usual finite measurable
partition gives this form by choosing its block index; null exceptional
points may be assigned to any block. The Lean predicate even omits the
measurability requirement on $b$, so its negation excludes a broader class
than the finite-step graphons in the source. The passage from a partition
written in ordinary graphon language to this index map, and the
identification of $[0,1]$ with its restricted-measure model, are semantic
translations explained here; no theorem concluding nonexistence is taken
as a Lean assumption.

The source's informal ``step count'' and ``jump count'' are not assigned
new definitions. The refuted existence assertion only needs the necessary
finite-range property proved in the file. The formal proofs contain no
unfinished terms or additional axioms: the printed dependency audit lists
only \texttt{propext}, \texttt{Classical.choice}, and \texttt{Quot.sound}.
No finite experiment or numerical approximation is used as evidence for
the infinite statement.

\paragraph{Source.}
\href{https://github.com/The-Last-Math-Competition/The-Last-Math-Competition/blob/main/conjectures/00000004408.md}{The Last Math Competition, Conjecture 00000004408}.
The exact bilingual source and build instructions are included in
\texttt{SOURCE.md} and \texttt{README.md}; snapshot identifiers are in
\path|verification/SOURCE_PROVENANCE.md|.

\end{document}
